\chapter{Curvature}
\label{chap:curvature}

% Closed question: What records the failure of transport to commute?

\begin{chapteressay}

Take a book. Rotate it 90 degrees east. Then rotate it 90 degrees
north. Observe the book's orientation. Now start with another copy
of the book in the same initial position. Rotate it 90 degrees north
first, then 90 degrees east. The two books are not in the same final
orientation. They differ by a 90-degree rotation about the vertical
axis.

The two sequences of rotations do not commute. The order matters. The
mathematical record of this non-commutativity is the chapter's
subject.

Curvature is the algebraic record of transport's failure to commute.
The previous chapter introduced transport: how information moves
between records. This chapter asks: when the transport is applied
in two different orders, does it produce the same result? The
answer is: not always. The cases where the result differs are
the cases of non-zero curvature. The chapter develops the
mathematics of curvature.

The chapter develops curvature through five structures. The first
is the commutator as the algebraic record of non-commutativity:
$[A, B] = AB - BA$. For rotation operators, the commutator is
the Lie bracket of the rotation algebra so(3). For parallel
transport on curved surfaces, the commutator is the Riemann
curvature tensor.

The second is Stokes' theorem as the integration of the
non-commutativity over a region. The boundary's circulation
detects the interior's curl; the closed-loop transport reveals
the enclosed curvature without requiring access to the
interior. The principle generalizes to all dimensions and all
admissible differential forms.

The third is the M\"obius band as a topological obstruction to
consistent transport. The band is locally flat but globally
non-trivial; parallel transport around the full band reflects
any vector. The chapter's lesson is that local flatness does
not imply global triviality; some non-trivial structure can be
topological rather than geometric.

The fourth is the three geometries: Euclidean (flat), spherical
(constant positive curvature), and hyperbolic (constant negative
curvature). The three are the canonical cases of constant
curvature. The Gauss-Bonnet theorem relates the integrated
curvature of a closed surface to its topological invariant
(the Euler characteristic).

The fifth is the heartbeat in Lean's elaboration as the
algorithmic analog of path-dependent strain. Different
elaboration paths accumulate different total heartbeats; the
difference is the chapter's algebraic record of the path's
elaboration strain.

The chapter's ground example is the rotation experiment. Rotating
a book east-then-north differs from rotating it north-then-east
by a $90^\circ$ rotation. The non-commutativity is observable
with a household object; the mathematical formalization is the
chapter's content.

The chapter's second concrete example is binding folded
signatures. A bindery has specific operations: fold, stack,
trim, glue, case-bind. The operations do not commute: changing
the order produces defective books. The binder's training has
internalized the admissible order; the discipline is the
chapter's calibration of the bindery's operations.

The chapter's closed question is:

\begin{center}
\emph{What records the failure of transport to commute?}
\end{center}

The chapter's answer is: the commutator algebra, formalized as
the Riemann curvature tensor for geometric transport, as the
boundary circulation in Stokes' theorem, as the topological
obstruction in the M\"obius band, and as the heartbeat
accumulation in Lean's elaboration. Each is a different
algebraic record of the non-commutativity.

Subsequent chapters move toward closure. Chapter 10 introduces
invariance: what survives admissible relabeling. The
non-commutativity of the current chapter is the algebraic
record of what differs; the invariance of the next chapter is
the record of what does not.

\end{chapteressay}

\section{Curvature as Commutator}
\label{se:curvature-as-commutator}

The essay above demonstrated the non-commutativity with a book and
two rotations: the east-then-north sequence ends in a different
orientation than the north-then-east sequence, differing by a
$90^\circ$ rotation about the vertical axis. This section moves
directly to the algebraic record.

The commutator of two operations is the difference between their
two compositions. For matrices $A$ and $B$,
\[
[A, B] = AB - BA.
\]
The commutator is zero when the operations commute and non-zero
otherwise. For the book's rotations, write $R_x$ for the rotation
about the east-west axis and $R_y$ for the rotation about the
north-south axis. The commutator $[R_x, R_y] = R_x R_y - R_y R_x$
is non-zero, and its specific value records the residual rotation
the two orderings differ by.

This is the chapter's first claim. The non-commutativity of
transport operations is the curvature. For rotations in
three-dimensional space, the rotation group SO(3) has non-zero
commutators; the commutators generate the Lie algebra so(3).
The non-commutativity is intrinsic to the geometry of rotation.

For parallel transport on a curved surface, the same phenomenon
appears. Transport a vector along path A then path B; transport
the same vector along path B then path A. The two results
differ. The difference is the holonomy of the closed loop
(A followed by reverse of B). The holonomy is proportional to
the surface curvature enclosed.

The Riemann curvature tensor $R^a_{bcd}$ is the mathematical
formalization of this non-commutativity. It is defined by the
commutator of covariant derivatives:
\[
[\nabla_c, \nabla_d] V^a = R^a_{bcd} V^b,
\]
where $V$ is a vector field and $\nabla$ is the covariant
derivative. The Riemann tensor measures the failure of parallel
transport to commute.

For flat spaces, the Riemann tensor is identically zero. Parallel
transport commutes; any closed loop accumulates zero holonomy.
For curved spaces, the Riemann tensor is non-zero somewhere.
The non-zero components quantify the curvature.

The chapter's claim is that curvature is the algebraic record of
non-commutativity. The mathematics is built from the order of
operations: which transport first, which second. The
commutator measures the dependence on order. The curvature is
the commutator's structural content.

\section{Stokes' Theorem as Yarn}
\label{se:stokes-theorem-as-yarn}

Stokes' Theorem in vector calculus relates the circulation of a
vector field around a closed loop to the flux of the field's curl
through any surface bounded by that loop:
\[
\oint_{\partial S} \mathbf{F} \cdot d\mathbf{r} = \int\int_S
(\nabla \times \mathbf{F}) \cdot d\mathbf{A}.
\]
The left side is a one-dimensional integral around the loop $\partial
S$. The right side is a two-dimensional integral over the surface
$S$ enclosed by the loop. The two are equal.

Stokes' Theorem is the chapter's second claim. The non-commutativity
of transport (the curl on the right) is detected by the
circulation around a closed loop (the integral on the left). The
loop's circulation measures the curl enclosed. The closed loop is
the chapter's diagnostic instrument: traversing it reveals the
curvature inside.

The theorem generalizes. In its most general form (the generalized
Stokes' theorem), it says: for any differential form $\omega$
on a manifold $M$ with boundary $\partial M$,
\[
\int_M d\omega = \int_{\partial M} \omega,
\]
where $d$ is the exterior derivative. The classical theorems
(divergence theorem, Green's theorem, the original Stokes'
theorem) are all special cases.

The principle is: boundary integrals detect interior structure.
The interior need not be explicitly traversed; the boundary's
behavior reveals what the interior contains. The chapter's
algebra carries this: the boundary's circulation is admissible
as a record of the interior's curvature.

The metaphor of yarn comes from the geometric picture. Imagine
threads laid across a surface, each one carrying the transport
of a vector. Around a closed loop, the threads accumulate a
specific twist; the twist is the holonomy. The twist comes from
the integrated curl along the threads' path. Yarn theory is the
algebraic study of these twists.

In Lean, the analog is the typeclass \texttt{YarnTheory} in
\texttt{Episode11.lean}:

\begin{leanexcerpt}{YarnTheory}
inductive YarnTheory
  | stokes : Fact $\to$ SpaceTimePath $\to$ Prop $\to$ YarnTheory
  | fibers : Fact $\to$ SpaceTimePath $\to$ SpaceTimePath
           $\to$ Prop $\to$ Prop $\to$ YarnTheory $\to$ YarnTheory
  | fabric : Fact $\to$ Fact $\to$ SpaceTimePath $\to$ SpaceTimePath
           $\to$ SpaceTimePath $\to$ Prop $\to$ Prop $\to$ Prop
           $\to$ YarnTheory $\to$ YarnTheory $\to$ YarnTheory
\end{leanexcerpt}

The inductive carries three layered shapes. The \texttt{stokes}
constructor sits at the base: a single \texttt{Fact}, a
\texttt{SpaceTimePath}, and a \texttt{Prop} that witnesses the
loop integral. The \texttt{fibers} constructor extends a prior
theory by attaching two paths and two propositional witnesses
under a \texttt{Fact}: the fiber assignments along the boundary.
The \texttt{fabric} constructor pairs two prior theories under a
triple of paths and a triple of propositional witnesses with two
\texttt{Fact}s: the curvature tensor as woven structure. Together
they specify the curvature's algebraic content.

The chapter's claim is that Stokes' theorem is the universal
detection principle. Curvature is interior structure; the
boundary's circulation is admissible as a record. Without
having to enter the interior, the closed-loop transport detects
the curvature. The discipline is the chapter's: the boundary
detects what the interior contains.

\section{The M\"obius Band}
\label{se:the-mobius-band}

Take a long strip of paper. Bring the two short ends together with
one end given a half-twist. Tape them. The result is a M\"obius
band: a surface with only one side and only one edge.

The M\"obius band is the simplest non-orientable surface. An ant
walking along the band's center will return to its starting
position after walking the full length of the band, but on the
opposite side --- a side that does not exist independently of the
original. The band's topology forbids a consistent ``up'' and
``down''.

Mathematically, the M\"obius band is a non-trivial line bundle
over the circle $S^1$. The line bundle is a structure assigning to
each point on the circle a one-dimensional fiber (a line); for the
M\"obius band, the fibers are connected in a twisted way as one
traverses the circle. Parallel transport of a vector in a fiber
around the full circle returns it reflected: a vector pointing
``up'' returns pointing ``down''.

The chapter's third claim is that the M\"obius band exhibits a
topological obstruction to consistent transport. The transport
around the loop must reflect any vector, because the bundle is
non-trivial. The reflection is the holonomy of the non-trivial
bundle.

The M\"obius band's curvature is not localized; it is global. The
band is locally flat: every small region of the band is
isometric to a strip of flat paper. But the band's global
topology is non-trivial; the obstruction is at the topological
level, not at the local geometric level.

The example illustrates the chapter's broader principle. Some
non-commutativity is local (the rotation example, the
Riemann curvature). Some is topological (the M\"obius band, the
fundamental group of a non-simply-connected space). The
chapter's algebra handles both: local curvature is recorded by
the Riemann tensor; topological obstructions are recorded by
the bundle's classifying class.

In Lean, the analog appears in the typeclass cascade for
non-trivial transport. A bundle's transport is admissible only
when the bundle is trivial or when the transport's holonomy is
explicitly accounted for. The
\texttt{YarnTheory.le} relation orders bundles by their
triviality: trivial bundles are at the bottom; non-trivial
bundles like the M\"obius band are higher in the order.

\section{The Three Geometries}
\label{se:the-three-geometries}

Euclidean geometry is the geometry of flat space. The parallel
postulate holds (given a line and a point not on it, there is
exactly one line through the point parallel to the given line).
Triangles have angle sums equal to $180^\circ$. The sum of the
squares of the legs of a right triangle equals the square of the
hypotenuse.

Spherical geometry is the geometry of a sphere. The parallel
postulate fails (any two ``lines'' --- great circles --- intersect).
Triangles have angle sums greater than $180^\circ$; the excess is
proportional to the triangle's area. The Pythagorean theorem
is modified.

Hyperbolic geometry is the geometry of the hyperbolic plane. The
parallel postulate fails in the opposite way (given a line and a
point, infinitely many lines through the point are parallel to
the given line). Triangles have angle sums less than $180^\circ$;
the deficit is proportional to the area.

These are the three geometries. They are distinguished by their
curvature: Euclidean has zero curvature; spherical has constant
positive curvature; hyperbolic has constant negative curvature.
Every two-dimensional homogeneous space with constant curvature
is locally one of the three.

The chapter's fourth claim is that the three geometries are the
canonical instances of curvature. The non-commutativity of
transport (Section 1) and the global obstruction (Section 3) are
both manifested in the three geometries. The Riemann tensor's
single non-zero component (in two dimensions) is the scalar
curvature; the scalar curvature is constant for these three
geometries.

For higher dimensions, the geometries multiply. The
$n$-dimensional sphere has constant positive curvature; the
$n$-dimensional hyperbolic space has constant negative curvature;
flat $n$-dimensional space has zero curvature. There are also
non-constant-curvature manifolds: surfaces of revolution, ellipsoids,
spacetimes in general relativity.

The Gauss-Bonnet theorem relates the total curvature of a
two-dimensional manifold to its topology. For a closed surface
$M$ with Euler characteristic $\chi(M)$:
\[
\int_M K \, dA = 2\pi \chi(M),
\]
where $K$ is the Gaussian curvature. For the sphere
($\chi = 2$), the integral is $4\pi$, recovering the standard
result that the total curvature of a sphere is $4\pi$. For a
torus ($\chi = 0$), the integral is $0$: some regions have
positive curvature, others negative, and they cancel exactly.

The chapter's claim is that curvature integrates to a
topological invariant. Local curvature is geometric; the integrated
curvature over a closed surface is topological. The connection
between geometry and topology is the chapter's deepest
structural claim.

\section{Heartbeat as Strain}
\label{se:heartbeat-as-strain}

A closed loop in a curved space accumulates a holonomy when a vector
is parallel-transported around it. The holonomy is the angle by
which the vector has rotated relative to its starting direction.
The rotation is the chapter's record of the enclosed curvature.

In Lean's elaboration, the analog is the heartbeat accumulation
along a proof path. A proof in Lean is constructed step by step;
each step consumes some elaboration cost (the heartbeat); the
total cost is the integrated heartbeat along the path.

If the proof is rewritten through a different sequence of tactics
(but arriving at the same conclusion), the heartbeat count differs.
The difference between the two paths' costs is the chapter's
elaborate heartbeat residue: the strain of having taken one path
rather than the other.

The chapter's fifth claim is that the heartbeat is the
algorithmic analog of curvature. Two elaboration paths
producing the same final result may accumulate different total
heartbeats. The difference is the path-dependent residue. The
algorithm's heartbeat-counting is the chapter's algebraic
record of the path's elaboration strain.

In Lean, the typeclass \texttt{HeartbeatProcess} formalizes this:

\begin{leanexcerpt}{HeartbeatProcess}
structure HeartbeatProcess
    (Value : Type)
    (Carrier : CarrierProcess Value)
    [DISTINGUISHABLE Value Carrier] ... [UNIVERSAL Value Carrier]
  where
  bullshit_meter       : CalculusProcess Value Carrier
  current_reading      : SpaceTimePath
  accumulated_bullshit : YarnTheory
  weave? : YarnTheory $\to$ YarnTheory := ...
\end{leanexcerpt}

The structure carries the project's full cascade through
\texttt{UNIVERSAL}, a \texttt{bullshit\_meter} (the
\texttt{CalculusProcess} that supplies the heartbeat reference), a
\texttt{current\_reading} along the spacetime path, an
\texttt{accumulated\_bullshit} as a yarn-theory record, and a
function \texttt{weave?} that advances the accumulated yarn by one
layer. As the proof elaborates, the structure updates the
accumulated yarn rather than a flat integer.

The compiler enforces the heartbeat cap. Each definition has a
maximum heartbeat budget (e.g., 400,000 by default). When the
budget is exceeded, the elaboration fails. The cap is the
chapter's resource constraint: the strain is bounded by the
budget; exceeding the budget means the path is too long.

The chapter's claim is that the heartbeat is the
metric on the proof's elaboration path. Different paths through
the proof state accumulate different heartbeats. The
heartbeat-difference between two paths is the elaboration
strain: the algebraic record of the path's curvature in the
proof space.

The chapter closes here. Curvature is the algebraic record of
non-commutativity. The Riemann tensor formalizes local
curvature; Stokes' theorem relates boundary circulation to
interior curvature; the M\"obius band shows topological
obstruction; the three geometries are the canonical
constant-curvature instances; the heartbeat is the
algorithmic analog of path-dependent strain.

\begin{bridge}

The chapter's question was what records the failure of transport to
commute.

The answer is curvature. The non-commutativity of transport
operations is the commutator; the commutator is the Riemann
tensor's structural content. Stokes' theorem relates boundary
circulation to interior curvature: the boundary detects the
interior. The M\"obius band illustrates topological
obstruction: even locally flat surfaces can have non-trivial
global topology. The three geometries (Euclidean, spherical,
hyperbolic) are the canonical constant-curvature instances.
The heartbeat in Lean's elaboration is the algorithmic
analog of path-dependent strain.

What remains is invariance. The chapter has identified the
non-commutativity; what about the structures that commute?
Chapter 10 takes up the question: what remains invariant
under admissible relabeling?

\end{bridge}

\begin{coda}{Binding Folded Signatures}

The bindery is at the back of the printshop, behind the press. The
binder has been at the trade for forty-two years. Her work today is
preparing the case binding for a small print run of a novel: five
hundred copies, each with three hundred sixty pages, hardcover with
a cloth-covered case.

The pages of the novel have been printed on large sheets, each sheet
containing sixteen pages: four pages on the front, four on the back,
with the proper layout so that when the sheet is folded, the pages
come out in order. The printed sheets are stacked beside the binding
table; there are twenty-three sheets per copy (totaling three hundred
sixty-eight pages, with eight blank pages at the back for end-papers
and overrun).

The binder takes a sheet and folds it. The fold is the first
operation: a single sharp crease running parallel to the short edge,
turning the sheet into a folded signature with eight pages visible
(four on each side after the fold). She folds again, this time
parallel to the long edge of the now-folded sheet: a second sharp
crease, producing a four-fold signature of sixteen pages.

The signatures of this novel are eight-page (one fold) for the
first and last sheets and sixteen-page (two folds) for the middle
sheets. The variation accommodates the novel's specific page count
and the binder's experience with what binds well.

After folding, the signatures are stacked in order. The order matters:
the first sheet's signature comes first, the second sheet's signature
second, and so on. If the signatures are stacked in the wrong order,
the pages will be out of order in the finished book; the order of
folding and the order of stacking must match.

After stacking, the signatures are trimmed. The folded edges (the
spine edges) are cut to a uniform depth; the top, bottom, and
foredge are cut to a uniform width. The trimming is done with a
guillotine cutter: a heavy blade descends with a controlled force
through the stacked signatures, producing a clean cut.

The sequence of operations is: fold, stack, trim. The chapter's
question is: does this sequence commute? What if the binder
trimmed before stacking, then folded? What if she stacked before
folding, then folded the whole stack?

The answer is: no, the operations do not commute. The folded
edges of the signatures are what hold the pages together in the
correct order. If the binder trimmed before stacking, the pages
within each signature might be out of register (the fold marks
would not align). If she stacked unfolded sheets and then folded
the whole stack, the inner pages would be longer than the outer
pages (because the fold radius increases for each layer), and
the bound book would have ragged pages.

This is the chapter's curvature in physical form. The bindery's
operations have a specific admissible order; the non-commutative
residue (the misregistered pages, the ragged edges) is what the
discipline avoids. The binder's training has internalized the
admissible order; the order is the chapter's calibration of
the bindery's operations.

After trimming, the spine of the stacked signatures is prepared
for binding. The binder applies glue to the spine edge,
inserts the spine into the case (the prepared cover), and
clamps the assembly to dry. The glue is the chapter's
\texttt{YarnTheory}: the fabric that holds the signatures
together through the case binding.

The chapter's heartbeat analog appears in the binder's pacing.
Each book takes a specific amount of time: the folding takes
about ten minutes for one copy; the trimming takes thirty
seconds; the spine gluing takes a minute; the drying takes an
hour. The total time per book is roughly fifteen minutes of
active work plus an hour of drying. The binder produces about
four books per hour during active assembly.

If the binder rushes, the operations may not commute correctly:
a misregistered fold here, a slightly-too-quick glue application
there. The strain of rushing accumulates as defects. The binder
slows down when her pace shows strain; she returns to the
admissible pace for the operations to commute correctly.

The five hundred copies will be bound over the next two weeks. The
binder will pace herself: roughly four books per hour, eight hours
per day, five days per week. The total work time is one hundred
twenty-five hours; the calendar time is sixteen working days. The
schedule is the chapter's path through the bindery's work: the
binder traverses the operations one signature at a time, one book
at a time, maintaining the admissible pace.

When the run is complete, the books will be inspected. Any with
misregistered folds, ragged edges, or weak spines will be
identified and returned for repair. The defects are the chapter's
curvature: the failures of the operations to commute correctly.
The bindery's record of repairs is the algebraic version of the
chapter's curvature accounting: each defect tells which operation
failed to commute at which step, and the repair restores the
admissible order.

\end{coda}
